\section{Check the complete program as a set of callable stages}
Keep the source project and a short verification record together. Record the input,
selected scalar or MSA mode, returned status and count, and the observation that
establishes the expected result. A displayed image is one observation; exact words,
unchanged padding and a failure before publication establish different properties.
Use the worked image first, then the fifteen-, sixteen- and seventeen-pixel inputs,
and finally the full-size repeated-row image.

Retain callable entries for loading, collection, inflation, reconstruction, packing
and complete decoding. Guided observations use the named construction boundaries;
a free-writing solution may replace their internal loops, scratch registers and
private frame layout. Keep the public arguments, saved machine state, counts and
memory ownership rules. A faster solution is accepted on the same output and access
contracts, rather than on similarity to an illustrated instruction sequence.

For each rejected input, identify the first stage that can establish the defect.
A damaged chunk CRC belongs to collection, an Adler mismatch to inflation, and a
validly compressed selector five to reconstruction. A short final framebuffer belongs
to packing admission. All four cases must leave the visible framebuffer and final
context fields unchanged. Scratch workspaces have a different failure contract: an
earlier inflater may have written an incomplete result before detecting its error.
Do not infer framebuffer damage merely from changed scratch bytes.

Include one success followed by a failure in the same session. Retain the successful
framebuffer and published context outputs before submitting the malformed image.
The failure must preserve those actual earlier results, not merely an initial zero
or sentinel pattern. Then decode a fresh valid image to show that stale table,
predictor or cursor state does not leak into the next call. Repeat in both modes.

For a performance observation, name the input, entry and return boundary, and count
simulated instructions over that interval. Compare exact output first. If a change
reduces work on the repeated Up rows but increases work on a small image, retain
both observations and explain the setup cost. The earlier prototype measurements
provide a reproducible comparison point; they are not target counts that every
correct implementation must reproduce.
